% Part: history
% Chapter: biographies
% Section: gerhard-gentzen

\documentclass[../../../include/open-logic-section]{subfiles}

\begin{document}

\olfileid{his}{bio}{gen}

\olsection{ګېرهارډ ګېنټسن}

\olphoto{gentzen-gerhard}{ګېرهارډ ګېنټسن}

ګېرهارډ ګېنټسن په عمده توګه د ثبوتونو د جوړښتي تيورۍ
د بنسټ ايښودونکي په توګه پېژندل کېږي، او په ځانګړي
ډول د طبيعي استنتاج او د سېکوېنټ حساب د
!!{derivation} سيسټمونو د جوړولو له امله۔ هغه د
1909 د نومبر پر 24مه د جرمني په ګرايفسوالډ کښې
زېږېدلے و۔ ګېرهارډ درې کاله په کور کښې زده کړه
وکړه او بيا د پوهنتون د چمتووالي ښوونځي ته لاړ،
چې هلته د زده کړې له مخې له خپلو ډېرو ټولګيوالو
وروسته و۔ له دې سره سره هغه تکړه زده کوونکے و
او په رياضياتو کښې ئې پياوړې وړتيا څرګنده کړه۔
شوقونه ئې بېلابېل وو؛ د بېلګې په توګه، خپلې مور
ته ئې شعرونه او د ښوونځي د تياتر دپاره ئې
ډرامې هم ليکلې۔

ګېنټسن خپلې پوهنتوني زده کړې د ګرايفسوالډ په
پوهنتون کښې پيل کړې، خو بيا ګوټينګن، ميونيخ او
برلين ته لاړ۔ په 1933 کښې ئې د ګوټينګن له
پوهنتون څخه د هېرمان وايل تر لارښوونې لاندې
دوکتورا واخيسته۔ (پاول برنايس ئې د کار د زياتې
برخې لارښوونه کړې وه، خو نازيانو هغه له پوهنتونه
ګوښه کړ۔) په 1934 کښې ګېنټسن د ډېوېډ هېلبرټ د
مرستيال په توګه کار پيل کړ۔ په هماغه کال ئې په
خپلو ليکنو \emph{Untersuchungen \"{u}ber das logische Schlie\ss en I--II
  [د منطقي استنتاج څېړنې I--II]} کښې د سېکوېنټ
حساب او طبيعي استنتاج !!{derivation} سيسټمونه
جوړ کړل۔ په 1936 کښې ئې د پيانو د بديهياتو
سازګاري ثابته کړه۔

له نازيانو سره د ګېنټسن اړيکه پېچلې ده۔ په
هماغه وخت کښې چې د هغه لارښوونکے برنايس د
جرمني پرېښودلو ته اړ کړے شو، ګېنټسن د نازي
نيمه‌پوځي سازمان SA په پوهنتوني څانګه کښې شامل
شو۔ د ډېرو جرمنيانو په شان هغه د نازي ګوند
غړے و۔ د جګړې پر وخت ئې د هوايي استخباراتو
د يوې څانګې د مخابراتي افسر په توګه کار کاوۀ۔
خو په 1942 کښې د عصبي ماتې له امله له دندې
خلاص کړے شو۔ دا څرګنده نۀ ده چې د ګېنټسن
وفاداري له نازي ګوند سره وه، که په ګوند کښې
د دې دپاره شامل شوے و چې علمي برياليتوب
ځان ته يقيني کړي۔

په 1943 کښې ګېنټسن ته د پراګ د جرمني پوهنتون
په رياضيکي انسټيټيوټ کښې علمي دنده وړاندې شوه،
چې هغه ومنله۔ خو په 1945 کښې د پراګ اوسېدونکو
د جرمني د اشغال پر ضد پاڅون وکړ۔ شوروي ځواکونه
ښار ته راورسېدل او د پوهنتون ټول پروفېسران ئې
ونيول۔ په نازي سازمانونو کښې د غړيتوب له امله
ګېنټسن د جبري کار يوه کېمپ ته بوتلل شو۔ هغه
په خپله حجره کښې د خوارځواکۍ له امله د 1945
د اګست پر 4مه د 35 کلونو په عمر وفات شو۔

\begin{reading}
د ګېنټسن د بشپړ ژوندليک دپاره
\citet{Menzler-Trott2007} وګورئ۔ د نازي واکمنۍ
لاندې د رياضيپوهانو په هکله يوه په زړه پورې
ليکنه، چې د ګېنټسن د ژوند لنډ يادښت هم ورکوي،
\citet{Segal2014} وړاندې کوي۔ د منطقي استنتاج
په هکله د ګېنټسن ليکنې په اصلي جرمني ژبه
شته \citep{Gentzen1935a,Gentzen1935b}۔ د ګېنټسن
د ليکنو انګرېزي ژباړې \citet{Gentzen1969} په
يو ټوک کښې راټولې کړې دي، چې د هغه د ژوند
يوه لنډه خاکه هم لري۔
\end{reading}

\end{document}
